\chapter{Eisenstein Series and the Continuous Spectrum}

\label{ch:8}

As before let $\bR^{\psi}$ denote the natural representation of $\bG_{\A}$ in $\bL^{2}(\bG_{\Q}\backslash\bG_{\A})$. The purpose of this section is to isolate the continuous spectrum of $\bR^{\psi}$ from the discrete spectrum. This will be accomplished by constructing an intertwining operator between a certain direct integral of induced representations and the restriction of $\bR^{\psi}$ to an appropriate invariant subspace whose orthocomplement decomposes discretely. Since this construction involves the theory of Eisenstein series and is not completely straightforward some additional remarks might be in order.

Our motivating theme here is the principle which asserts that a unitary representation can be decomposed by spectrally decomposing the self-adjoint operators which commute with it. In the context of $\bR^{\psi}(g)$ this means (as in Section~\ref{ch:2}) that functions in $\bL^{2}(\bG_{\Q}\backslash\bG_{\A},\psi)$ must be expanded in terms of $\psi$-automorphic forms on $\GL(2)$. Eisenstein series enter the picture because they produce sufficiently many automorphic forms to decompose the ``continuous'' part of $\bL^{2}(\bG_{\Q}\backslash\bG_{\A},\psi)$. Obstacles arise because these forms themselves are not square-integrable and therefore an intermediary Fourier transform is required. The reader might keep in mind that this situation is analogous to what happens when the regular representation of $\R$ is decomposed using the (non square-integrable) ``forms'' $e^{ixy}$.

We must emphasize that it is not our purpose to develop the theory of Eisenstein series from scratch. Rather our purpose is to explain the role Eisenstein series play in the decomposition of $\bR^{\psi}(g)$ and (in Section~\ref{ch:9}) in the proof of the Selberg trace formula. Therefore several facts from the theory of Eisenstein series will be recalled without proof. For details and complete proofs the reader is referred to \cite{Selberg}, \cite{Langlands4}, \cite{HarishChandra2}, \cite{Langlands5}, \cite{Kubota} and ``Further Notes and References'' at the end of this section.

\section{Some Preliminaries}\label{sec:8a}

Let $Z_{\infty}^{+}$ denote the group

\[
\left\{\begin{bmatrix} r & 0 \\ 0 & r \end{bmatrix} \in Z_{\infty}: r > 0\right\}
\]

and put

\begin{equation}\label{eq:8.1}
\bX = Z_{\infty}^{+}\bG_{\Q}\backslash\bG_{\A}.\tag{8.1}
\end{equation}

Since it is slightly more convenient (but no less general) to deal with $\bL^{2}(\bX)$ in place of $\bL^{2}(\bG_{\Q}\backslash\bG_{\A},\psi)$ we shall do so until further notice. Therefore $\bR(g)$ will henceforth denote the natural (unitary) representation of $\bG_{\A}$ in $\bL^{2}(\bX)$. In Subsections C and D the theory of Eisenstein series will be used to produce the preliminary decomposition

\begin{equation}\label{eq:8.2}
\bR = \bR_{0}^{+}\oplus\bR_{1}\tag{8.2}
\end{equation}

where $\bR_{0}^{+}$ is a direct sum of certain (unknown) irreducible representations and $\bR_{1}$ is a direct integral (or continuous sum) of certain (known) induced representations $\bR(g:z)$, $z \in i\R$.

Throughout this section the following notation will be fixed: $B$ will denote the upper triangular group of $\bG$, $A$ the diagonal subgroup, and $N$ the upper niltriangular subgroup. The Haar measure $\dd x$ will denote the Tamagawa measure on $\bG_{\A}$. This choice of measure is made for normalization purposes only so its precise definition is not important. See [Ono] or \cite{Weil6} for details.

On any discrete subgroup of $\bG_{\A}$ we shall use the Haar measure which assigns to any point the measure one, and on any quotient of unimodular groups to which we have assigned Haar measures we shall use the corresponding quotient measure. Thus if we assign to $Z_{\infty}^{+}$ the measure which corresponds to Euclidean measure on $(\R_{+})^{\times}$ we determine a measure on $\bG_{\A}^{1} = Z_{\infty}^{+}\backslash\bG_{\A}$, and hence also a measure on

\[
\bX = Z_{\infty}^{+}\bG_{\Q}\backslash\bG_{\A}
\]

whose total volume is the Tamagawa number of $\GL(2)$.

Similarly, we have uniquely determined measures on $A_{\A}$, $N_{\A}$, $A_{\infty}^{+} = \left\{\begin{bmatrix} a_{1} & 0 \\ 0 & a_{2} \end{bmatrix} \in A_{\infty}: a_{i} > 0\right\}$, and

\begin{equation}\label{eq:8.3}
A_{\A}^{1} = \left\{\begin{bmatrix} a_{1} & 0 \\ 0 & a_{2} \end{bmatrix} \in A_{\A}: |a_{1}| = |a_{2}| = 1\right\} = A_{\infty}^{+}\backslash A_{\A}.\tag{8.3}
\end{equation}

On $B_{\A}$, which is non-unimodular, we define a left Haar measure by

\begin{equation}\label{eq:8.4}
\int_{B_{\A}} \phi(b) \dd_{\ell} b = \int_{A_{\A}} \int_{N_{\A}} \phi(an) \dd a \dd n\tag{8.4}
\end{equation}

and right Haar measure by

\begin{equation}\label{eq:8.5}
\int_{B_{\A}} \phi(b) \dd_{r} b = \int_{N_{\A}} \int_{A_{\A}} \phi(na) \dd n \dd a.\tag{8.5}
\end{equation}

We recall that

\begin{equation}\label{eq:8.6}
\dd_{r} b = \delta_{B}(b) \dd_{\ell} b\tag{8.6}
\end{equation}

where

\[
\delta_{B}(b) = \left|\frac{a_{1}}{a_{2}}\right|, \quad \text{if} \quad b = \begin{bmatrix} a_{1} & n \\ 0 & a_{2} \end{bmatrix}.
\]

Also that $\dd(ana^{-1}) = \delta_{B}(a) \dd n$ and that

\begin{equation}\label{eq:8.7}
\int_{\bG_{\A}} f(x) \dd x = c_{G} \int_{B_{\A}} \int_{\bK} f(bk) \delta_{B}(b)^{-1} \dd_{r} b \dd k\tag{8.7}
\end{equation}

for some constant $c_{G}$.

Now fix once and for all an isomorphism

\[
t \to h_{t} = \begin{bmatrix} e^{t} & 0 \\ 0 & e^{-t} \end{bmatrix}
\]

from $\R$ onto a subgroup $T$ of $A_{\infty}^{+}$ (so $A_{\infty}^{+}$ is the direct product of $Z_{\infty}^{+}$ and $T$). Write

\begin{equation}\label{eq:8.8}
x = na h_{t} k z\tag{8.8}
\end{equation}

where $n \in N_{\A}$, $a \in A_{\A}^{1}$, $t \in \R$, $k \in \bK$, and $z \in Z_{\infty}^{+}$. Since the number $t$ is uniquely determined by $x$ we may denote it by $H(x)$. Note that

\begin{equation}\label{eq:8.9}
H(b) = \log \left|\frac{a_{1}}{a_{2}}\right|^{1/2} \quad \text{if} \quad b = \begin{bmatrix} a_{1} & 0 \\ 0 & a_{2} \end{bmatrix},\tag{8.9}
\end{equation}

and

\begin{equation}\label{eq:8.10}
\delta_{B}(b) = e^{2H(b)} \quad \text{if} \quad b \in B_{\A}.\tag{8.10}
\end{equation}

The integration formula corresponding to the decomposition \eqref{eq:8.8} is then

\begin{equation}\label{eq:8.11}
\int_{Z_{\infty}^{+}\backslash \bG_{\A}} f(x) \dd x = c_{G} \int_{N_{\A}} \int_{A_{\A}^{1}} \int_{-\infty}^{\infty} \int_{\bK} f(na h_{t} k) e^{-2t} \dd n \dd a \dd t \dd k.\tag{8.11}
\end{equation}

Finally we recall the notion of Siegel domain appropriate to this context. For $c > 0$ put

\begin{equation}\label{eq:8.12}
\gamma(c) = \omega A_{\infty}^{+}(c)\bK\tag{8.12}
\end{equation}

where $\omega$ is some relatively compact subset of $N_{\A}A_{\A}^{1}$ and

\begin{equation}\label{eq:8.13}
A_{\infty}(c) = \{a \in A_{\infty}^{+}: H(a) \geq \log c\}.\tag{8.13}
\end{equation}

The Siegel domain defined by \eqref{eq:8.12} may be chosen so that $\bG_{\A} = \bG_{\Q}\gamma(c)$.

A function $f(x)$ on $\bG_{\A}$ will be called slowly increasing on $\gamma(c)$ if there are constants $C$ and $N$ such that

\[
|f(x)| \leq Ce^{-NH(x)}, \quad x \in \gamma(c),
\]

and rapidly decreasing on $\gamma(c)$ if for every $N$ there is a constant $C_{N}$ such that

\[
|f(x)| \leq C_{N}e^{NH(x)}, \quad x \in \gamma(c).
\]

If $f$ is any continuous function on $Z_{\infty}^{+}\bG_{\Q}\backslash\bG_{\A}$ its constant term is the function

\begin{equation}\label{eq:8.14}
f^{0}(x) = \int_{N_{\Q}\backslash N_{\A}} f(nx) \dd n.\tag{8.14}
\end{equation}

Note that $f$ is a cusp form if and only if its constant term is zero almost everywhere. If $f$ is an automorphic form on $\GL(2)$ (in the sense of Definition~\ref{def:3.3}, with $\psi$ trivial on $Z_{\infty}^{+}$) then $f - f^{0}$ is rapidly decreasing on any Siegel domain. In particular $f - f^{0}$ is square integrable on $\bX$. Consequently every cuspidal automorphic form is in $\bL^{2}(\bX)$, as already remarked in Section~\ref{ch:3}.

\section{Analysis of Certain Induced Representations}\label{sec:8b}

The representations in question are induced from the (non-unimodular) subgroup $N_{\A}A_{\Q}A_{\infty}^{+}$ of $B_{\A}$. For each $z \in \C$ consider the character

\[
b = \begin{bmatrix} a_{1} & u \\ 0 & a_{2} \end{bmatrix} \to e^{zH(b)} = \left|\frac{a_{1}}{a_{2}}\right|^{z/2}
\]

of $N_{\A}A_{\Q}A_{\infty}^{+}$. The resulting induced representation

\[
\bR(x,z) = \underset{N_{\A}A_{\Q}A_{\infty}^{+} \uparrow \bG_{\A}}{\Ind} e^{zH(b)}
\]

operates by right translation in the space $\bH(z)$ consisting of complex-valued measurable functions $\phi$ on $\bG_{\A}$ such that

\begin{equation}\label{eq:8.15}
\int_{\bK} |\phi(k)|^{2} \dd k < \infty\tag{8.15}
\end{equation}

and

\begin{equation}\label{eq:8.16}
\phi(nag) = e^{(z+1)H(na)}\phi(g) = \left|\frac{a_{1}}{a_{2}}\right|^{1/2 + z/2}\phi(g)\tag{8.16}
\end{equation}

for all $na \in N_{\A}A_{\Q}A_{\infty}^{+}$. Note that $\bR(x,z)$ is actually a representation of $Z_{\infty}^{+}\backslash\bG_{\A}$.

\begin{lemma}\label{lem:8.1}

Let $\bH$ denote the Hilbert space of measurable functions $\phi$ on $N_{\A}A_{\Q}A_{\infty}^{+}\backslash\bG_{\A}$ satisfying \eqref{eq:8.15} and define a representation $\bR'(x,z)$ in $\bH$ through the formula

\begin{equation}\label{eq:8.17}
\bR'(x,z)\phi(g) = \phi(gx)e^{(z+1)H(gx)}e^{-(z+1)H(g)}.\tag{8.17}
\end{equation}

Then the representations $\bR(x,z)$ and $\bR'(x,z)$ are equivalent. The operator intertwining them is

\begin{equation}\label{eq:8.18}
I(z): \phi(g) \to \phi(g)e^{-(z+1)H(g)}.\tag{8.18}
\end{equation}

\end{lemma}

\begin{proof}

Completely straightforward since $H(nag) = H(na) + H(g)$.

\end{proof}

\begin{lemma}\label{lem:8.2}

The Hilbert space adjoint of $\bR(x,z)$ is $\bR(x^{-1}, -\overline{z})$. Consequently $\bR(x,z)$ is unitary whenever $\operatorname{Re}(z) = 0$.

\end{lemma}

\begin{proof}

To prove this lemma it is convenient to use the ``multiplier realization'' of $\bR(x,z)$, namely $\bR'(x,z)$. Indeed since $N_{\A}A_{\Q}A_{\infty}^{+}\backslash\bG_{\A}$ is compact we can find a real-valued function $\beta$ in $C_{c}^{\infty}(\bG_{\A})$ such that $\int_{N_{\A}A_{\Q}A_{\infty}^{+}}\beta(bx)\dd_{r}b = 1$ for $x \in \bG_{\A}$. Consequently

\begin{equation}\label{eq:8.19}
(\phi, \psi) = \int_{\bG_{\A}} e^{2H(x)}\beta(x)\phi(x)\overline{\psi(x)} \dd x\tag{8.19}
\end{equation}

for $\phi, \psi$ in $\bH$ (by \eqref{eq:8.7} and \eqref{eq:8.15}). The desired equality

\[
(\bR'(x,z)\phi, \psi) = (\phi, \bR'(x^{-1}, -\overline{z})\psi)
\]

then follows from \eqref{eq:8.19} by an obvious change of variables.

\end{proof}

Our aim is to show that the direct integral

\begin{equation}\label{eq:8.20}
\int_{\substack{\operatorname{Re}(z) = 0 \\ \operatorname{Im}(z) > 0}} \bR(x,z) \dd|z|\tag{8.20}
\end{equation}

(of unitary representations $\bR(x,z)$) is equivalent to a subrepresentation of (the unitary) representation $\bR(x)$. We shall not need the representations $\bR(x,z)$ with negative imaginary part since $\bR(x,-z)$ will turn out to be equivalent to $\bR(x,z)$.

Before explaining why \eqref{eq:8.20} is natural we shall need to further analyze its summands. Indeed the representations $\bR(x,z)$ are themselves not irreducible and it is of some interest to know how they decompose.

By Mackey's ``Inducing in Stages Theorem'' $\bR(x,z)$ may be described by first inducing $e^{zH(na)}$ from $N_{\A}A_{\Q}A_{\infty}^{+}$ to $N_{\A}A_{\A}$ and then inducing the resulting representation (call it $T(z)$) from $N_{\A}A_{\A}$ to $\bG_{\A}$. Now the representation space of $T(z)$ consists of functions on $N_{\A}A_{\A}$ which are (at least) left $N_{\A}$-invariant. Therefore $T(z)$ may be viewed as a representation of $A_{\A}$. In fact if $\bR_{A}$ denotes the regular representation of $A_{\Q}\backslash A_{\A}^{1}$ (a compact abelian group) then $T(z)$ is equivalent to the tensor product of $\bR_{A}$ with the one-dimensional representation $e^{(z+1)H(a)}$ of $A_{\infty}^{+}$. (Since $H\!\left(\begin{bmatrix} a_{1} & 0 \\ 0 & a_{2} \end{bmatrix}\begin{bmatrix} a_{1}' & 0 \\ 0 & a_{2}' \end{bmatrix}\right) = H\!\left(\begin{bmatrix} a_{1} & 0 \\ 0 & a_{2} \end{bmatrix}\right) + H\!\left(\begin{bmatrix} a_{1}' & 0 \\ 0 & a_{2}' \end{bmatrix}\right)$ the action of $T(z)$ on $\bL^{2}(A_{\Q}A_{\infty}^{+}\backslash A_{\A})$ is given by $\phi(y) \to \phi(yx)e^{(z+1)H(x)}$.) Consequently the following proposition is immediate.

\begin{proposition}\label{prop:8.3}

Let $\Lambda$ denote the set of characters $\lambda$ of $A_{\Q}A_{\infty}^{+}\backslash A_{\A}$. Let $\bR(x,\lambda,z)$ denote the principal series representation of $\bG_{\A}$ induced from the character $e^{zH(a)}\lambda(a)$ of $N_{\A}A_{\A}$. Then

\begin{equation}\label{eq:8.21}
\bR(x,z) = \bigoplus_{\lambda \in \Lambda} \bR(x,\lambda,z).\tag{8.21}
\end{equation}

\end{proposition}

The definition below makes it possible to complete our analysis of the space of $\bR(x,z)$.

\begin{definition}\label{def:8.4}

Suppose $\tau$ is an irreducible representation of $\bK$. Then we say that $\phi \in C^{\infty}(N_{\A}A_{\Q}A_{\infty}^{+}\backslash\bG_{\A})$ is of type $\tau$ if for any $x \in \bG_{\A}$, $\phi(xk)$ is contained in a subspace of $\bL^{2}(\bK)$ on which the right regular representation of $\bK$ is equivalent to $\tau$ (i.e.\ $\phi(xk)$ transforms according to $\tau$).

\end{definition}

The space of all such $\phi$ will be denoted by $(\tau)$; this space is finite-dimensional and has a natural inner product given by

\[
(\phi_{1}, \phi_{2}) = c_{G}\int_{\bK} \phi_{1}(k)\overline{\phi_{2}(k)} \dd k.
\]

If $L$ denotes the collection of all such $(\tau)$ then

\begin{equation}\label{eq:8.22}
\bH = \bigoplus_{(\tau) \in L} (\tau).\tag{8.22}
\end{equation}

\section{Eisenstein Series}\label{sec:8c}

Roughly speaking, Eisenstein series provide a map from the space of $\bR(g,z)$ to a space of automorphic forms on $\bG_{\A}$. More precisely this theory provides an isomorphism between the space of \eqref{eq:8.20} and the continuous spectrum of $\bL^{2}(\bX)$ which commutes with the natural action of $\bG_{\A}$. The basic Eisenstein series are defined as follows.

For each $\phi \in (\tau)$, $(\tau) \in L$ (cf.\ Definition~\ref{def:8.4}) consider the function

\begin{equation}\label{eq:8.23}
(I^{-1}(z)\phi)(g) = \phi(x)e^{(z+1)H(g)}.\tag{8.23}
\end{equation}

Since $I^{-1}(z)$ is the ``inverse'' of $I(z)$ this function belongs to $\bH(z)$ (since $\tau \in \bH$). In particular it transforms according to \eqref{eq:8.16} under the action of $N_{\A}A_{\Q}A_{\infty}^{+}$. The Eisenstein series associated to $\phi$ is (by definition) the series

\begin{equation}\label{eq:8.24}
\begin{aligned}
E(g,\phi,z) &= \sum_{\gamma \in B_{\Q}\backslash\bG_{\Q}} (I^{-1}(z)\phi)(\gamma g)\\
&= \sum_{\gamma \in B_{\Q}\backslash\bG_{\Q}} \phi(\gamma g)e^{(z+1)H(\gamma g)}.
\end{aligned}\tag{8.24}
\end{equation}

\begin{remark}\label{rem:8.5}

The Eisenstein series just introduced generalizes the ``classical'' Eisenstein series

\begin{equation}\label{eq:8.25}
\begin{aligned}
E(x + iy, z) &= \sum_{\gamma \in \Gamma \cap N \backslash \Gamma} f_{\mu}(\gamma(x + iy))\\
&= \sum_{\substack{(c,d) = 1 \\ c,d \in \Z}} \frac{y^{\frac{z+1}{2}}}{|cz + d|^{z+1}}
\end{aligned}\tag{8.25}
\end{equation}

described in Example~\ref{exa:2.5}. Indeed $E(g,\phi,z) = E(x + iy,z)$ precisely when $\phi$ is identically one and $g = \begin{bmatrix} y & x \\ 0 & 1 \end{bmatrix}$ is real, $y > 0$.

\end{remark}

In general it is known from the theory of Selberg-Langlands that

(1) the series defining $E(x,\phi,z)$ converges uniformly for $x$ in compact subsets of $Z_{\infty}^{+}\backslash\bG_{\A}$ and $z$ in compact subsets of

\begin{equation}\label{eq:8.26}
D = \{z: \operatorname{Re}(z) > 1\};\tag{8.26}
\end{equation}

(2) for $z \in D$ this series defines a $\bK$-finite eigenfunction of $\mathfrak{z}$ (the center of the universal enveloping algebra of $\bG_{\infty}$) which is slowly increasing in all Siegel domains, i.e.\ $E(x,\phi,z)$ is an automorphic form;

(3) the constant term of $E(x,\phi,z)$, i.e.

\begin{equation}\label{eq:8.27}
E^{0}(x,\phi,z) = \int_{N_{\Q}\backslash N_{\A}} E(nx,\phi,z) \dd n\tag{8.27}
\end{equation}

is given by the expression

\begin{equation}\label{eq:8.28}
\phi(x)e^{(z+1)H(x)} + M(z)\phi(x)e^{(-z+1)H(x)}\tag{8.28}
\end{equation}

where $M(z)$ is a uniquely determined analytic function (depending on $(\tau)$) which maps $D$ into the space of linear operators on $(\tau)$: the adjoint of $M(z)$ is $M(\overline{z})$;

(4) $E(x,\phi,z)$ and $M(z)\phi$ can be analytically continued as meromorphic functions of $z$ regular on $i\R$: the only possible pole to the right of $i\R$ is simple and occurs at $z = 1$, simultaneously for $E(x,\phi,z)$ and $M(z)\phi$;

(5) $E(x,\phi,z)$ is slowly increasing if $\operatorname{Re}(z) = 0$; and

(6) the following functional equations are satisfied:

\begin{equation}\label{eq:8.29}
M(z)M(-z)\phi = \phi,\tag{8.29}
\end{equation}

and

\begin{equation}\label{eq:8.30}
E(x,M(z)\phi,-z) = E(x,\phi,z).\tag{8.30}
\end{equation}

We note that if $\tau$ is trivial then $M(z)$ is equal to $\xi(z)/\xi(z+1)$ with

\[
\xi(z) = \pi^{-z/2}\Gamma(z/2)\zeta(z).
\]

Therefore the analytic continuation and functional equation of $M(z)$ is (in this case) quite believable.

The significance of the facts (1)-(6) for the decomposition of $\bR$ can be explained as follows. When $\bR(z) = 0$, both $\bR(x,z)$ and $E(x,\phi,z)$ are defined (the latter in the sense of analytic continuation). The representation $\bR(x,z)$ is unitary by Lemma~\ref{lem:8.2}, and $E(x,\phi,z)$ is an automorphic form by (2) and (5) above. Our claim is that the map

\begin{equation}\label{eq:8.31}
\phi \to E(x,\phi,z)\tag{8.31}
\end{equation}

``essentially'' intertwines the space of $\bR(x,z)$ with a subrepresentation of $\bR(x)$.

Indeed a (formal) computation shows that

\begin{equation}\label{eq:8.32}
E(x,\bR(g,z)\phi,z) = \bR(g)E(x,\phi,z)\tag{8.32}
\end{equation}

i.e. that the map \eqref{eq:8.31} from $\bH$ to some space of automorphic forms on $\bG_{\A}$ commutes with $\bG_{\A}$. The problem is that $E(x,\phi,z)$ is not square integrable and therefore (as remarked earlier) a more intricate Fourier transform is required. Before describing this transform let us directly construct a certain continuous sum of the ``principal series'' representations $\bR(x,z)$.

For each $(\tau) \in L$ let $\hat{\bL}_{0}^{2}(E,\tau)$ denote the (one or zero dimensional) space of functions from $\{1\}$ to $(\tau)$ with positive semi-definite inner product

\begin{equation}\label{eq:8.33}
\frac{1}{2\pi}\left(a_{0}(1), \mu(1)b_{0}(1)\right).\tag{8.33}
\end{equation}

Here $\mu(1)$ is the residue of $2\pi i M(z)$ at $z = 1$ and functions are identified modulo null vectors.

Let $\hat{\bL}_{1}^{2}(E,\tau)$ denote the space of square-integrable functions $a_{1}(z)$ from $i\R$ to $(\tau)$ such that

\begin{equation}\label{eq:8.34}
a_{1}(-z) = M(z)a_{1}(z)\tag{8.34}
\end{equation}

and define an inner product in this space

\begin{equation}\label{eq:8.35}
(a_{1}, b_{1}) = \frac{1}{\pi}\int_{\operatorname{Re}(z) = 0} (a_{1}(z), b_{1}(z)) \dd|z|.\tag{8.35}
\end{equation}

Note that the Hilbert space

\[
\bigoplus_{(\tau) \in L} \hat{\bL}_{1}^{2}(E,(\tau))
\]

is essentially the direct integral of the Hilbert spaces $\bH(z)$ with respect to Lebesgue measure on $\{z = iy, y > 0\}$. Finally put

\begin{equation}\label{eq:8.36}
\hat{\bL}^{2}(E,\tau) = \hat{\bL}_{0}^{2}(E,\tau) \oplus \hat{\bL}_{1}^{2}(E,\tau).\tag{8.36}
\end{equation}

Now let $\bH(\tau)$ denote the space of entire functions with values in $(\tau)$ which are Fourier-Laplace transforms of functions in $C_{c}^{\infty}(i\R) \otimes (\tau)$.

For any $a \in \bH(\tau)$ define

\begin{equation}\label{eq:8.37}
\tilde{a}(x) = \frac{1}{2\pi}\int_{\operatorname{Re}(z) = z_{0}} E(x,a(z),z) \dd|z|\tag{8.37}
\end{equation}

where $z_{0}$ is arbitrary but to the right of 1. Again from Selberg-Langlands' theory it is known that

(7) $\tilde{a}(x) \in \bL^{2}(\bX)$;

(8) if the closure of $\{\tilde{a}: a \in \bH(\tau)\}$ in $\bL^{2}(\bX)$ is denoted by $\bL^{2}(E,\tau)$ then

\begin{equation}\label{eq:8.38}
\bL^{2}(E) = \bigoplus_{(\tau) \in L} \bL^{2}(E,\tau)\tag{8.38}
\end{equation}

is precisely the orthocomplement (in $\bL^{2}(\bX)$) of the space of cusp forms $\bL_{0}^{2}(\bX)$; and

(9) for any $a, b$ in $\bH(\tau)$,

\begin{equation}\label{eq:8.39}
\begin{aligned}
\int_{\bX} \tilde{a}(x)\tilde{b}(x) \dd x
&= \frac{1}{2\pi}\int_{\operatorname{Re}(z) = z_{0}} \left\{(a(z), b(-\overline{z})) + (M(z)a(z), b(\overline{z}))\right\} \dd|z|\\
&= \frac{1}{2\pi}(\mu(1)a(1), b(1))\\
&\quad + \frac{1}{4\pi}\int_{-i\infty}^{i\infty} (a(z) + M(-z)a(-z), b(z) + M(-z)b(-z) \dd|z|.
\end{aligned}\tag{8.39}
\end{equation}

This last equality results from an application of the residue theorem (which is justified by the fact that the norm of $M(z)$ is bounded at infinity in the strip $\{0 \leq \operatorname{Re}(z) \leq z_{0}\}$.

We shall now explain the significance of these facts for the decomposition of $\bR(x)$. Using (7)-(9) one can define an isomorphism $\bE$ between $\hat{\bL}^{2}(E,\tau)$ and $\bL^{2}(E,(\tau))$ (and hence ultimately between $\hat{\bL}^{2}(E)$ and a certain privileged subspace of $\bL^{2}(\bX)$) by requiring that $\bE$ map $(a_{0},a_{1})$ to $\tilde{a}$ whenever $(a_{0},a_{1})$ arises from $a \in \bH(\tau)$ through the formulas

\[
a_{0}(1) = a(1)
\]

and

\begin{equation}\label{eq:8.40}
a_{1}(z) = \frac{1}{2}(a(z) + M(-z)a(-z)) \quad (z \in i\R).\tag{8.40}
\end{equation}

Indeed by (9) the correspondence

\[
(a_{0}, a_{1}) \xleftrightarrow{\bE} \tilde{a}
\]

is a linear isometry between dense subspaces of $\hat{\bL}^{2}(E,\tau)$ and $\bL^{2}(E,\tau)$. On the other hand by (8) the orthocomplement of $\oplus\bL^{2}(E,(\tau))$ is $\bL_{0}^{2}(\bX)$ and by Theorem~\ref{thm:5.1} this latter space decomposes discretely. Consequently the role Eisenstein series play in the decomposition of $\bR(x)$ is clear.

\section{Description of the Continuous Spectrum}\label{sec:8d}

Recall that the goal of this section was to construct an invariant subspace of $\bL^{2}(\bX)$ whose orthocomplement exhausts the discrete spectrum of $\bR(x)$. Accepting the theory of Eisenstein series this goal is near at hand.

Our claim is that properties (7)-(9) of the last subsection infer that the subspace of $\bL^{2}(\bX)$ defined by

\[
\bigoplus_{(\tau) \in L} \bE(\hat{\bL}_{1}^{2}(E,\tau))
\]

is this sought after subspace.

To see this, let $\bE_{0}$ (respectively $\bE_{1}$) denote the restriction of the isomorphism $\bE$ to $\hat{\bL}_{0}^{2}(E,\tau)$ (resp.\ $\hat{\bL}_{1}^{2}(E,\tau)$). By taking direct sums over all $(\tau) \in L$ we obtain maps from

\[
\hat{\bL}_{i}^{2}(E) = \bigoplus_{(\tau) \in L} \hat{\bL}_{i}^{2}(E,\tau) \quad (i = 1,2)
\]

which we again denote by $\bE_{i}$. If we denote by $\bL_{i}^{2}(E)$ the image of $\bE_{i}$ then by (7) each $\bL_{i}^{2}(E)$ belongs to $\bL^{2}(\bX)$ and by (8) and (9)

\begin{equation}\label{eq:8.41}
\bL^{2}(\bX) = \bL_{0}^{2}(\bX) \oplus \bL_{0}^{2}(E) \oplus \bL_{1}^{2}(E).\tag{8.41}
\end{equation}

For each compactly supported measurable function $f$ on $Z_{\infty}^{+}\backslash\bG_{\A}$ we now introduce the (bounded) operators

\begin{equation}\label{eq:8.42}
\bR(f,z) = \int_{Z_{\infty}^{+}\backslash\bG_{\A}} f(g)\bR(g,z) \dd g\tag{8.42}
\end{equation}

and

\begin{equation}\label{eq:8.43}
\bR(f) = \int_{Z_{\infty}^{+}\backslash\bG_{\A}} f(g)\bR(g) \dd g.\tag{8.43}
\end{equation}

With this notation the intertwining property of Eisenstein series (originally expressed by \eqref{eq:8.32}) takes the form

\begin{equation}\label{eq:8.44}
\bR(f)E(x,\phi,z) = E(x,\bR(f,z)\phi,z).\tag{8.44}
\end{equation}

This last identity is an identity between meromorphic functions of $z$ and it is meaningful in a rigorous sense as soon as $f$ is bi-$\bK$-finite (since $\bR(f,z)\phi$ belongs to a finite collection of $(\tau)$'s whereas $\bR(x,z)\phi$ need not).

Now the spectrum of $\bL_{0}^{2}(\bX)$ is known to be discrete by Theorem~\ref{thm:5.1} so to prove our claim it remains to observe that $\bL_{0}^{2}(E)$ too is invariant with discrete spectrum. This claim in turn will follow from the fact that $\hat{\bL}_{0}^{2}(E)$ owes its existence to poles of Eisenstein series to the right of $i\R$ and such poles occur only at $z = 1$.

In fact such poles occur at $z = 1$ only when $\tau(k) = \chi(\det k)$ with $\chi$ a character of $\Q^{\times}\backslash\A^{\times}$ whose square is trivial on $\R_{+}^{\times}$. The corresponding residue of $E(x,\phi,z)$ is then proportional to the function

\[
\bD_{\chi}(g) = \chi(\det g).
\]

Thus a consequence of the lemma below is that $\bL_{0}^{2}(E)$ is equivalent to the direct sum of the one-dimensional representations $\bD_{\chi}(g)I$. This lemma gives an explicit formula for the functions $\bE_{i}a_{i}(x)$ and is crucial to the sequel. We include its proof because the computations involved are so typical of the general theory.

\begin{lemma}\label{lem:8.6}

Suppose $(a_{0},a_{1})$ is such that $a_{0} \in \hat{\bL}_{0}^{2}(E,\tau)$ and $a_{1}$ is a function of compact support in $\hat{\bL}_{1}^{2}(E,\tau)$. Then for almost all $x$,

\begin{equation}\label{eq:8.45}
\begin{aligned}
\bE_{0}a_{0}(x) + \bE_{1}a_{1}(x) &= i\operatorname{Res}_{z=1} E(x,a_{0}(1),z)\\
&\quad + \frac{1}{2\pi}\int_{-i\infty}^{i\infty} E(x,a_{1}(z),z) \dd|z|.
\end{aligned}\tag{8.45}
\end{equation}

\end{lemma}

\begin{proof}

Let $f_{0}(x) + f_{1}(x)$ denote the right side of \eqref{eq:8.46}. To prove the lemma it will suffice to prove that the function $g(x) = f_{0}(x) + f_{1}(x) - \bE_{0}a_{0}(x) - \bE_{1}a_{1}(x)$ satisfies the following properties:

\begin{enumerate}[(i)]

\item it is square-integrable;

\item it is orthogonal to $\bL_{0}^{2}(\bX)$; and

\item its constant term vanishes almost everywhere.

\end{enumerate}

Indeed these properties imply $g(x) = 0$ a.e.

For convenience, let $E_{0}(x,\phi,z)$ denote $2\pi i\operatorname{Res}E(x,\phi,z)$. Then $E_{0}$ itself is an automorphic form and consequently $E_{0} - E_{0}^{0}$ is rapidly decreasing on any Siegel domain $\gamma(c)$. But from \eqref{eq:8.27} and \eqref{eq:8.28} it follows that $E_{0}^{0}(x,a_{0}(1),1) = (\mu(1)a_{0}(1)(x)$. Thus

\[
\begin{aligned}
\int_{\gamma(c)} |E_{0}^{0}(x,a_{0}(1),1)|^{2} \dd x &= \int_{\omega}\int_{\log c}^{\infty}\int_{\bK} |(\mu(1)a_{0}(1))(na h_{t} k)|^{2}e^{-2t} \dd b \dd t \dd k\\
&= \left(\int_{\log c}^{\infty} e^{-2t} \dd t\right)\int_{\omega}\int_{\bK} (\mu(1)a_{0})(1)(ak) \dd b \dd k < \infty,
\end{aligned}
\]

i.e.\ $E_{0}^{0}$ is square-integrable. So since $E_{0} - E_{0}^{0}$ is rapidly decreasing, $E_{0}(x,a_{0}(1),1)$ (and hence $f_{0}(x)$) is square integrable too.

Now consider the term $f_{1}(x)$. Since $a_{1}(z)$ has compact support, and $E(x,\phi,z)$ is an automorphic form,

\[
|f_{1}(x) - f_{1}^{0}(x)| \leq \int_{-i\infty}^{i\infty} \left|E_{1}(x,a_{1}(z),z) - E_{1}^{0}(x,a_{1}(z),z)\right| \dd|z| \leq C_{N}e^{-NH(x)}
\]

(for all $N$). Therefore $f_{1} - f_{1}^{0}$ is square integrable. But from (4.28) and the fact that $M(z)a_{1}(z) = a_{1}(-z)$ it is easy to see that

\begin{equation}\label{eq:8.46}
f_{1}^{0}(x) = \frac{1}{\pi}\int_{-i\infty}^{i\infty} a_{1}(z)(x)e^{(z+1)H(x)} \dd|z|.\tag{8.46}
\end{equation}

This function, being the Fourier transform of the rapidly decreasing function $a_{1}(z)$, is also square integrable, and therefore $f_{1}(x)$ and $g(x)$ are square integrable, as desired.

To prove (ii) it suffices to deal with $f_{0}(x) + f_{1}(x)$ alone since $(\bE_{0}a_{0})(x)$ and $(\bE_{1}a_{1})(x)$ are already known to be orthogonal to $\bL_{0}^{2}(\bX)$. So suppose $h$ is an arbitrary automorphic form in $\bL_{0}^{2}(\bX)$. Then $h$ is rapidly decreasing, and for any $\phi \in (\tau)$, $z \in D = \{\operatorname{Re}(z) > 1\}$,

\[
\begin{aligned}
\int_{\bX} E(x,\phi,z)\overline{h(x)} \dd x &= \int_{Z_{\infty}^{+}\bG_{\Q}\backslash\bG_{\A}} \sum_{\gamma \in B_{\Q}\backslash\bG_{\Q}} \phi(\gamma x)e^{(z+1)H(\gamma x)}\overline{h(\gamma x)} \dd x\\
&= \int_{Z_{\infty}^{+}B_{\Q}\backslash\bG_{\A}} (\phi(x)e^{(z+1)H(x)})\overline{h(x)} \dd x\\
&= \int_{A_{\Q}\backslash A_{\A}^{1}}\int_{\bK} \phi(ak)e^{(z+1)H(a)}\left(\int_{N_{\Q}\backslash N_{\A}} h(nak) \dd n\right) \dd a \dd k\\
&= 0,
\end{aligned}
\]

since $h$ is cuspidal. That $f_{1}(x)$ and $f_{2}(x)$ are also orthogonal to $h$ follows by analytic continuation. So since $\bL_{0}^{2}(\bX)$ has a basis consisting of automorphic forms, the proof of (ii) is complete.

To prove (iii), choose a sequence of functions $\{b^{n}\}$ in $\bH(\tau)$ so that $(b_{0}^{n},b_{1}^{n})$ approaches $(a_{0},a_{1})$ in $\hat{\bL}^{2}(E,\tau)$. By the residue theorem, and the relation (4.28), the constant term of $\tilde{b}^{n}$ is

\[
\begin{aligned}
(b^{n})^{0}(x) &= \int_{N_{\Q}\backslash N_{\A}} \frac{1}{2\pi}\int_{\operatorname{Re}(z) = z_{0}} E(nx,b^{n}(z),z) \dd|z| \dd n\\
&= \frac{1}{2\pi}\int_{\operatorname{Re}(z) = z_{0}} \left\{b^{n}(z)(x)e^{(z+1)H(x)} + (M(z)b^{n}(z))(x)e^{(-z+1)H(x)}\right\} \dd|z|\\
&= \frac{1}{2\pi}(\mu(1)b_{0}^{n}(1))(x) + \frac{1}{\pi}\int_{\operatorname{Re}(z) = 0} (b_{1}^{n}(z))(x)e^{(z+1)H(x)} \dd|z|.
\end{aligned}
\]

But from \eqref{eq:8.46} it is then clear that by a suitable choice of $\{b^{n}\}$ this last expression can be made to approach $f_{0}^{0}(x) + f_{1}^{0}(x)$ for all $x$, as $n$ approaches $\infty$. Therefore the constant terms of $\{\bE_{0}a_{0} + \bE_{1}a_{1}\}$ and $\{f_{0} + f_{1}\}$ are equal almost everywhere, since $\tilde{b}^{n}$ approaches $\bE_{0}a_{0} + \bE_{1}a_{1}$ in the mean.

\end{proof}

\begin{remark}\label{rem:8.7}

From the crucial identity \eqref{eq:8.44} and the identity \eqref{eq:8.30} (the functional equation for $E(x,\phi,z)$) it follows that

\begin{equation}\label{eq:8.47}
\bR(f,z)M(z) = M(z)\bR(f,-z),\tag{8.47}
\end{equation}

i.e. the operator intertwining the equivalent representations $\bR(x,z)$ and $\bR(x,-z)$ is precisely $M(z)$ (summed over all $(\tau) \in L$).

\end{remark}

An immediate consequence of Lemma~\ref{lem:8.6} and the identity \eqref{eq:8.47} is the invariance of the subspace $\bL_{1}^{2}(E)$ and its equivalence (via $\bE_{1}$) to the direct integral \eqref{eq:8.20}.

\section{Summing Up}\label{sec:8e}

\begin{theorem}\label{thm:8.8}

Let $\bR_{0}$ denote the restriction of $\bR$ to the space of cusp forms and $\bR^{+}$ the restriction to $\bigoplus_{\chi}\C \bD_{\chi}$. (Recall that $\chi$ is any grossencharacter of $\Q$ whose square is trivial on $(\R_{+})^{\times}$ and $\bD_{\chi}(g) = \chi(\det g)$.) Then

\[
\bR = \bR_{0}^{+}\oplus\bR_{1}
\]

where

\[
\bR_{0}^{+}(x) = \bR_{0}(x)\oplus\bR^{+}(x)
\]

and

\[
\bR_{1}(x) = \int_{0}^{\infty} \bR(x,it) \dd t.
\]

Furthermore the projection of $\bL^{2}(\bX)$ onto the space of $\bR_{1}$ is $\bE_{1}\bE_{1}^{*}$.

\end{theorem}

\begin{corollary}\label{cor:8.9}

Theorem~\ref{thm:8.8} describes the continuous spectrum of $\bL^{2}(\Gamma\backslash\SL(1,\R))$ for congruence subgroups of $\SL(2,\Z)$.

\end{corollary}

\begin{example}[\normalfont (Theorem~\ref{thm:2.6})]\label{exa:8.10}

Let $\bH(s)$ denote the representation space of the principal series representation of $\SL(2,\R)$ induced from the character

\[
\begin{bmatrix} a & u \\ 0 & a^{-1} \end{bmatrix} \to |a|^{is}.
\]

Then if $\Gamma = \SL(2,\Z)$,

\[
\bL^{2}(\Gamma\backslash\SL(2,\R)) = \bL_{0}^{2}(\Gamma\backslash\SL(2,\R)) \oplus \C \oplus \int_{0}^{\infty} \bH(s) \dd s
\]

where $\C$ denotes the subspace of $\bL^{2}$ spanned by the constant functions.

\end{example}

\begin{proof}

As an $\SL(2,\R)$-module, $\bL^{2}(\Gamma\backslash\SL(2,\R))$ is isomorphic to the subspace of $\bL^{2}(Z_{\infty}^{+}\bG_{\Q}\backslash\bG_{\A})$ consisting of right $\bK_{0}$-invariant functions. So by Proposition~\ref{prop:8.3} (and the definition of $\bD_{\chi}(g)$) this example is an immediate consequence of Theorem~\ref{thm:8.8}.

\end{proof}

We close this section with a lemma which will be used in Section~\ref{ch:9}.

\begin{lemma}\label{lem:8.11}

Suppose $(\tau) \in L$ and $\phi \in (\tau)$. Then for any $h \in C_{c}^{\infty}(Z_{\infty}^{+}\bG_{\Q}\backslash\bG_{\A})$

\[
(\bE_{1}^{*}h(z),\phi) = \frac{1}{2}\int_{\bX} h(x)\overline{E(x,\phi,z)} \dd x
\]

for almost all $z \in i\R$.

\end{lemma}

\begin{proof}

Suppose $a_{1}(z)$ is an arbitrary continuous compactly supported function in $\hat{\bL}_{1}^{2}(E,\tau)$. Then

\[
\frac{1}{\pi}\int_{-i\infty}^{i\infty} ((\bE_{1}^{*}h)(z),a_{1}(z)) \dd|z| = \int_{\bX} h(x)\overline{\bE_{1}a_{1}(x)} \dd x
\]

since both sides equal the inner product of $\bE_{1}^{*}h$ and $a_{1}$ (the left side by definition of the inner product in $\hat{\bL}_{1}^{2}(E,\tau)$, the right side by definition of the adjoint of $\bE_{1}$). Interchanging the order of integration on the right side (justified since $a_{1}(z)$ is compactly supported) yield the expression

\[
\frac{1}{2\pi}\int_{-i\infty}^{i\infty} \left(\int_{\bX} h(x)\overline{E(x,a_{1}(z),z)} \dd x\right) \dd|z|.
\]

So by the arbitrariness of $a_{1}$ the desired conclusion follows.

\end{proof}

\section*{Notes and References}

\addcontentsline{toc}{subsection}{Notes and References}

The results from the theory of Eisenstein series sketched in this section are due to \cite{Selberg}. Unfortunately Selberg published them without giving complete proofs. The corresponding results for general semi-simple algebraic groups were proved by Langlands and recorded (with complete details) in \cite{Langlands4,Langlands5}. Similar (but weaker) results were independently obtained in \cite{GelfandGraevPyatetskiiShapiro}.

In addition to the references already given in the text the reader should consult \cite{HarishChandra3}. There the theory of Eisenstein series is sketched for reductive groups defined over a finite field.

Our discussion in this section is adapted from \cite{Arthur}. Similar summaries are to be found in \cite{DufloLabesse} and \cite{Godement2,Godement3}. We have worked over the rationals for the sake of simplicity but it goes without saying that the theory is valid for arbitrary global fields as well.

We note finally that the operator $M(z)$ of \eqref{eq:8.47} can be written as a product of local operators $M_{p}(z)$. Then the operator $M_{\infty}(z)$ essentially coincides with the intertwining operator of \cite{KunzeStein} and \cite{KnappStein}.
